\subsection{Анализ и выбор стека технологий}
В области компьютерного зрения, существует множество фреймворков, которые позволяют разработчикам быстро и эффективно создавать приложения для работы с видеопотоком в реальном времени. Фреймворки для анализа видеопотока в реальном времени обладают различными инструментами, такими как построение конвейеров, модули ввода и вывода, микширование видеопотока и подключение модулей нейросетей для детекции, классификации, сегментации и трекинга объектов. Все это помогает извлечь полезную информацию из видеопотока. 

Существует несколько популярных фреймворков, которые решают данные задачи: OpenCV, GStreamer\cite{Gstreamer}, Nvidia Deepstream \cite{Deepstream} и Savant. Рассмотрим все из них и выберем тот, который лучше подходит для решения наших задач.

OpenCV (Open Source Computer Vision Library) - это библиотека компьютерного зрения с открытым исходным кодом, которая предоставляет набор функций для обработки изображений и видео. Она была создана в 1999 году в Intel и с тех пор получила широкую популярность в научных кругах, индустрии и сообществе разработчиков.

OpenCV написана на языке программирования C++ и имеет биндинги для других языков, таких как Python, Java и MATLAB. Библиотека предоставляет реализацию многих алгоритмов компьютерного зрения, таких как сегментация изображений, детектирование объектов, распознавание лиц, оптическое распознавание символов (OCR), стереозрение и многое другое.

OpenCV также предоставляет инструменты для обработки видео, такие как захват видео, воспроизведение и запись видео, а также алгоритмы для трекинга объектов и анализа движения.

Недостаток OpenCV - это то, что некоторые его функции могут быть довольно медленными, особенно при работе с большими наборами данных или высокой разрешающей способностью изображений. Это может быть проблемой при работе с видеопотоками в реальном времени. Также, в OpenCV нет инструментов для обработки многих видеопотоков одновременно, что может быть проблематично при работе с большим количеством камер или источников видео. 

Таким образом, хотя OpenCV может использоваться для обработки видео в реальном времени, его функциональность ограничена, и для создания более сложных конвейеров обработки видео может потребоваться использование дополнительных инструментов и библиотек.


GStreamer - это открытая библиотека для обработки мультимедиа, которая предоставляет набор инструментов для создания и манипулирования потоками данных мультимедиа, такими как аудио и видео.

Она предоставляет множество компонентов, таких как кодеки, фильтры, элементы и плагины, которые могут быть использованы для создания цепочек обработки мультимедиа, которые могут обрабатывать потоки данных из разных источников, например, с веб-камер, микрофонов, файлов или сетевых потоков.

Один из ключевых преимуществ GStreamer заключается в ее гибкости и настраиваемости. С помощью GStreamer, разработчики могут создавать сложные цепочки обработки мультимедиа с использованием различных компонентов, фильтров и плагинов, которые могут быть настроены для удовлетворения конкретных потребностей.

GStreamer также имеет множество расширений и плагинов, которые позволяют разработчикам добавлять дополнительную функциональность в свои приложения. Например, есть плагины для обработки изображений, обработки звука, а также для работы с различными видео-кодеками и аудио-кодеками.

Функциональность обработки изображений с помощью нейронных сетей не является частью стандартного набора элементов GStreamer и не встроена в GStreamer непосредственно. Для интеграции нейронных сетей для обработки изображений необходимо использовать сторонние плагины или создать свой собственный плагин.


NVIDIA DeepStream это набор инструментов и библиотек, разработанных NVIDIA, который предоставляет разработчикам возможность создавать и развертывать приложения компьютерного зрения и аналитики для обработки видео в реальном времени на базе GPU. DeepStream является расширением GStreamer, которое добавляет поддержку инференса нейронных сетей, аналитики и очередей, что позволяет создавать приложения для обработки видео в реальном времени с использованием глубокого обучения. 

DeepStream обеспечивает высокую производительность и эффективность обработки видео благодаря использованию аппаратного ускорения на GPU и оптимизированных алгоритмов для работы с видео. С помощью DeepStream можно создавать приложения для решения различных задач, таких как детекция объектов, распознавание лиц, анализ поведения, подсчет людей, анализ трафика и т.д. DeepStream поддерживает многокамерную обработку видео и распределенную обработку данных на нескольких устройствах.

Кроме того, DeepStream интегрируется с другими инструментами NVIDIA, такими как TensorRT и CUDA, что позволяет разработчикам использовать широкий набор оптимизированных алгоритмов машинного обучения для обработки видео.

В целом, DeepStream является мощным инструментом для создания высокопроизводительных приложений компьютерного зрения и аналитики для обработки видео в реальном времени на базе NVIDIA GPU.

Savant - это технология, основанная на Nvidia DeepStream, которая предоставляет возможность работы с видеокадрами на оборудовании Nvidia с высокой производительностью. Она поддерживает платы Jetson, а также позволяет получить доступ к видеокадрам прямо из кода Python.

Однако, в отличие от Deepstream, Savant имеет меньшую гибкость и меньшее количество возможностей. Savant конвейер настраивается с помощью файла YAML, в который разработчик помещает цепочку блоков обработки. Эти блоки могут включать блоки, связанные с DeepStream, пользовательские блоки на основе Python и предварительно разработанные блоки Savant. В таблице \ref{tab:tab1} представлено сравнение технологий для работы с видео.

\begin{table}    
    \caption{Сравнительная таблица технологий для работы с видео}
    \begin{tabularx}{\textwidth}{|X|c|c|c|c|}\hline
         & OpenCV & GStreamer  & Nvidia Deepstream & Savant \\ \hline
    Поддержка построения конвейеров  & -   & +    & +   & +         \\ \hline
    Работа с нейросетями       & +   & -    & +  & +         \\ \hline
    Расширяемость   & $\pm$   & +    & +  & $\pm$         \\ \hline
    Поддержка вывода модели из сторонних фреймворков & -  & +   & +  & -         \\ \hline
    \end{tabularx}
    \label{tab:tab1}
\end{table}

%добавить сравнительную таблицу

Таким образом, если необходима более гибкая и функциональная технология для работы с видеокадрами, то Deepstream может быть более подходящим вариантом. Однако, если нужна более простая технология, которая может обеспечить высокую производительность на оборудовании Nvidia, то Savant может быть хорошим выбором.


Для быстрой работы моделей необходимо использовать их с помощью сервера вывода, о котором было описано в предыдущем пункте. Также рассмотрим уже существующие решения.

TensorFlow Serving - это система, разработанная Google для обслуживания моделей машинного обучения, построенных с помощью TensorFlow. Она позволяет развертывать обученные модели TensorFlow в производственной среде, что позволяет обрабатывать запросы на предсказания по сетевому интерфейсу с использованием различных API и протоколов.

TensorFlow Serving поддерживает как онлайн-предсказание, так и предсказание пакетами, что означает, что он может обрабатывать запросы на предсказания в режиме реального времени, а также запросы на обработку больших наборов данных в пакетном режиме. Он также включает функции, такие как версионирование моделей, гибкая загрузка моделей и устойчивость к сбоям.

Основная идея TensorFlow Serving состоит в том, чтобы разделить процесс обучения модели от ее использования в производственной среде. Обученная модель может быть экспортирована в формате, подходящем для TensorFlow Serving, и развернута на сервере, который будет обслуживать запросы на предсказания. Сервер TensorFlow Serving может использовать различные API для предоставления предсказаний, такие как gRPC, REST API или TensorFlow API.

ONNX Runtime - это кросс-платформенный движок глубокого обучения, который разработан Microsoft для эффективного выполнения моделей машинного обучения на различных устройствах и платформах. ONNX - это открытый формат обмена моделями машинного обучения, который был разработан совместно с Facebook и другими компаниями.

ONNX Runtime поддерживает различные форматы моделей машинного обучения, такие как TensorFlow, PyTorch, Keras, Scikit-learn, и другие после конвертации в ONNX формат. Это позволяет разработчикам машинного обучения создавать модели в любом из этих фреймворков и использовать ONNX Runtime для выполнения моделей на различных устройствах и платформах, таких как CPU, GPU, FPGA и мобильные устройства.

ONNX Runtime также имеет ряд оптимизаций для ускорения выполнения моделей, включая автоматическое использование аппаратного ускорения, оптимизации для уменьшения памяти и быстрого выполнения операций.

NVIDIA Triton - это набор инструментов и библиотек, разработанный NVIDIA для развертывания моделей машинного обучения на производственных системах. Он предоставляет высокую производительность и масштабируемость, что позволяет эффективно запускать модели на больших кластерах серверов.

Основным компонентом NVIDIA Triton является сервер развертывания моделей, который позволяет загружать и запускать модели на производственных системах. Он поддерживает различные форматы моделей, включая TensorFlow, PyTorch, ONNX и другие, и предоставляет простой интерфейс для управления моделями.

NVIDIA Triton также включает в себя ряд инструментов для управления моделями, мониторинга и настройки производительности. Например, он может автоматически оптимизировать работу моделей, чтобы максимизировать использование ресурсов GPU и снизить задержки при обработке запросов. Triton также поддерживает многопоточность и многопроцессорность, что позволяет обрабатывать большое количество запросов одновременно. Сравнение серверов вывода приведены в таблице  \ref{tab:tab2}.

%добавить сравнительную таблицу

\begin{table}    
    \caption{Сравнительная таблица серверов вывода}
    \begin{tabularx}{\textwidth}{|X|c|c|c|}\hline
         & Nvidia Triton & TensorFlow Serving  & ONNX Runtime\\ \hline
    Tensor RT  & +   & -    & -            \\ \hline
    ONNX       & +   & -    & +         \\ \hline
    TensorFlow   & +   & +    & -          \\ \hline
    Pytorch & +   & -    & -           \\ \hline
    Python Code & +   & -    & -   \\ \hline
    \end{tabularx}
    \label{tab:tab2}
\end{table}

Так как NVIDIA Triton обладаем большой производительностью на видеокартах Nvidia, поддерживает множество форматов нейронных сетей, имеет инструменты для построения конвейера моделей как единой модели и хорошей интеграцией с Nvidia Deepstream остановимся именно на этом решении.

В связи с выбранным набором технологий следуют следующие системные требования, приведенные в таблице \ref{tab:tab3}.

\begin{table}    
    \caption{Системные требования}
    \begin{tabularx}{\textwidth}{|X|c|}\hline
    Платформа & x86\_64, Jetson\\ \hline
    Ubuntu  & 20.04           \\ \hline
    Python       & 3.8        \\ \hline
    CUDA   & 11.4 для Jetson и 11.8 для x86\_64         \\ \hline
    GPU & RTX2080/RTX3080/T4/A100   \\ \hline
    DeepStream SDK & 6.2   \\ \hline
    Docker & >=19.03.3  \\ \hline
    \end{tabularx}
    \label{tab:tab3}
\end{table}